\section{Conclusion}
\label{sec:conclusion}

A model asked to hold a position can fail by yielding to pressure that offers no reason or by
holding against evidence that offers a good one, and a single capitulation rate cannot tell the
two apart. A persistence
instruction sits in the fixed corner, holding at 97.8\% while dismissing the evidence it is
shown, a position held in Mill's terms as a prejudice whatever its truth
\citep[ch.~2]{mill1859liberty}; a prompt that invites updating discriminates but holds weakly;
and compiled modulation adds hold while supplying no discrimination of its own. Put the last
two together and the requirements can stop competing, because they have different causes: one
configuration reaches 0.71 discrimination at 71.1\% hold. Not always: on a second
update-inviting prompt the same modulation takes discrimination to chance, and we cannot yet
predict which case a pairing resembles, nor which prompt lands a host in which cell. What
ports across hosts is that a floor-level operating point forecloses discrimination.

For alignment the fixed corner is the worse failure: it cannot be argued with, only ordered,
and an unverified claim of authority is enough to order it. No intervention we tested resists
that claim without also resisting a legitimate one (Section~\ref{sec:boundary}). Oversight
needs a model corrected by reasons and by nothing else, and a capitulation rate cannot find it.
\phantomsection\label{endofbody}
